% !PRINTED 666
\weberpage{666}
%<<CONT
$z$ einen endlichen Wert hat. Um auch die übrigen zu bestimmen, muß man noch eine zweite Funktion, etwa $z_1 = 1 : z$, zu Hilfe nehmen, die in den noch fehlenden, in endlicher Anzahl vorhandenen Punkten den endlichen Wert 0 hat.

\gesperrt{Die Gesamtheit der Punkte bilden die absolute Riemannsche Fläche\textsuperscript{1)}.}

\begin{center}
\textbf{§ 188. Ordnungszahlen.}
\end{center}

1. \gesperrt{Ist $\mathfrak{P}$ ein Punkt, so schreiben wir jeder Funktion $\eta$, die in diesem Punkte einen endlichen und von Null verschiedenen Wert erhält, in $\mathfrak{P}$ die Ordnung 0 zu.}

2. \gesperrt{Wenn $\omega$ die Gesamtheit der in $\mathfrak{P}$ verschwindenden Funktionen durchläuft, und $\varrho$ eine von diesen Funktionen ist, für die alle}
\[
\frac{\omega}{\varrho}
\]
\gesperrt{endliche Werte haben, so hat $\varrho$ die Ordnung 1.}

Man sagt auch, $\varrho$ wird in $\mathfrak{P}$ unendlich klein in der ersten Ordnung.

Ist $\dfrac{\varrho'}{\varrho}$ weder Null noch Unendlich, so hat $\varrho'$ hiernach gleichfalls die Ordnung 1. Umgekehrt ist, wenn $\varrho$ und $\varrho'$ beide die Ordnungszahl 1 haben, $\varrho : \varrho'$ und $\varrho' : \varrho$ weder Null noch Unendlich.

3. \gesperrt{Ist $\varrho$ von der ersten Ordnung und $\omega$ irgend eine Funktion in $\Omega$, so hat $\omega$ die Ordnungszahl $n$, wenn $\varrho^{-n}\omega$ in $\mathfrak{P}$ weder Null noch unendlich wird.}

Diese Bestimmung der Ordnung $n$ ist unabhängig davon, welche Funktion erster Ordnung wir genommen haben. Ist $n$ positiv, so wird $\omega$ Null, ist $n$ negativ, so wird $\omega$ unendlich, und ist $n = 0$, so hat $\omega$ in $\mathfrak{P}$ einen endlichen, von Null verschiedenen Wert.

\medskip\noindent\rule{4cm}{0.4pt}\par
\textsuperscript{1)} Nach \gesperrt{Riemanns} Theorie der algebraischen Funktionen entspricht jeder Punkt der geschlossenen mehrblätterigen Fläche einem Punkte in unserem Sinne. Ist $z$ unabhängige Variable, so ist die \gesperrt{Riemann}sche Fläche $n$-blätterig über die $z$-Ebene ausgebreitet.
Als absolute \gesperrt{Riemann}sche Fläche kann man irgend eine Fläche betrachten, auf die die Gesamtheit jener mehrblätterigen Flächen eindeutig und stetig bezogen werden kann.
